\chapter*{はじめに}
\addcontentsline{toc}{chapter}{はじめに}

このテキストは、Python を使ってデータを読み、整形し、図を描き、分布を調べるための入門書である。単に文法を覚えるのではなく、結果の正しさを自分で確かめ、処理時間を測り、改善点を説明できるようになることを目標にする。

本書の構成は Python 公式チュートリアルの流れを意識している。前半では、変数、リスト、関数、ファイル読み書きといった基礎を確認する。中盤では NumPy、pandas、Matplotlib、SciPy、Astropy を用いた実践的な解析へ進む。後半では、大規模データの高速化や低レベル言語との連携に触れ、最後に総合演習へ進む。

特に、Python 公式チュートリアルの「An Informal Introduction to Python」「More Control Flow Tools」「Data Structures」「Modules」「Input and Output」「Errors and Exceptions」「Brief Tour of the Standard Library」で扱われている基礎事項を、データ解析の文脈でつなぎ直すことを意識している。公式チュートリアルを置き換えるというより、その内容を土台にして、解析実装で本当に困る点まで掘り下げることを目指す。

コマンド例は、macOS 上のターミナルで作業する状況を自然な基準として書いている。
Terminal.app でも iTerm2 でも構わない。
Windows ユーザーは PowerShell を用いてもよいが、Windows Subsystem for Linux (WSL) の使用を強く推奨する。
